\chapter{Discussion}\label{ch:discussion}
This chapter discusses the design, capabilities and limitations of \benchmark, building on the concept, implementation and experiments presented in the previous chapters. Section \ref{sec:discussion-benchmark-design} discusses key design decisions behind \benchmark, focusing on its hierarchical scoring system and its configurable components. Section \ref{sec:evaluation-benchmark-principles} evaluates \benchmark against the principles of a well-designed benchmark introduced in Section \ref{ssec:principles-benchmark}. Section \ref{sec:existing-benchmark-comparison} compares \benchmark to the related benchmarking tools discussed earlier in this thesis. Section \ref{sec:experiment-learnings} reflects on the key learnings drawn from the experiments conducted in Chapter \ref{ch:experiments}, regarding how \ac{rl} agents for power grid operation should be benchmarked in general. Finally, Section \ref{sec:discussion-limitations} discusses the limitations of this work.




\section{Discussion of Benchmark Design}\label{sec:discussion-benchmark-design}
\benchmark offers a customisable benchmark that allows users to create experiments by giving them a selection of scenarios, with each scenario offering a chosen power grid and profile from the IEEE test cases and SimBench. The grid modifications offer the possibility to modify the grid data to create more diverse scenarios. This customisability introduces several design decisions that are discussed in this section. First, the hierarchical scoring system is discussed, covering how the underlying metrics are used, how they are aggregated across multiple layers, and how the resulting scores are presented back to the user. Second, the configurable components of the benchmark are discussed, namely the scoring weights and the scenario selection, which together define the objective and quality of a given benchmark trial.

\subsection{Scoring System}
The scoring system used by \benchmark structures the evaluation of an agent's performance across multiple layers, ranging from individual grid metrics to a single, aggregated total score. This subsection discusses the design of this scoring system, covering how the underlying metrics are used and normalised, how the resulting scores are aggregated across the hierarchy, and how the different layers are presented back to the user for interpretation.

\subsubsection{Metrics Usage}
As demonstrated in Section \ref{sec:metrics-eval}, each experiment has a hierarchical scoring system. On the lowest layer, there are multiple grid metrics that are measured during a single run. Multiple metrics were chosen, as there is not one single metric that can provide all the information needed on the agent's performance. Each of them focusses on only one grid or operational feature, covering multiple features that are important in evaluation of an agent in power grid operation. While most features focus on grid stability, some focus on operational cost or computational efficiency. All of these categories are relevant in a strong performance in power grid operation, as managing a power grid must also remain under a specific budget limit, without causing costs to rise. Also, the computational efficiency is crucial in real-time power grid operation, when the agent needs to react and decide on the next actions quickly, or when the computational resources are limited.

Furthermore, the metrics are converted into a normalised score from 0 to 100 by using a linear function, relative to the worst and best reference values. The normalisation into 0 to 100 makes every metric comparable and simple to interpret, as the score values are in a percentage-based range. Additionally, the normalisation allows aggregation across the metrics, due to the same value range for each metric. 

\subsubsection{Score Aggregation}
The first aggregation happens with categories, which group the metrics and provide a summary of how well the agent performed for specific types of grid metrics in one run. Each category covers a crucial aspect of strong performance in power grid operation.

Each run then receives a weighted average score of the categories, which is then used to calculate the average scenario score. This provides another summarised feedback for the next two layers, run and scenario layers. Without needing to see all the measured metrics and their scores, it is possible to directly assess the performance of each run and scenario. Finally, the highest level of summary is the total score, which is calculated as the average of the scenario scores. This gives feedback for the overall performance, without showing any details about the exact metric, category or other scores. 

Each layer in this hierarchical scoring system has its importance. The higher layers summarise information to provide a simple comparison value for quick evaluation. The lower layers gather specific data points and create scores for them, thus having the most information about individual performances. The lower layers are crucial for analysing the agent's behaviour and finding patterns or problems.

\subsubsection{Result Presentation}
Building on the multiple aggregated scores across different layers of the hierarchical scoring system, \benchmark offers many ways to present the results of a benchmark trial. As shown in Section \ref{sec:result-format}, the result presentation ranges from statistics, plots, web UIs for analysis and comparison, and pages for replay of a run. These presentations again emphasise the premise that one single score is not sufficient for meaningful feedback. They showcase each scenario and run, creating different textual reports and visualisations for them. Additionally, the comparison views allow contrasting multiple agents directly against each other on the same scenario or category, making it easier to identify in which specific aspects one agent outperforms another, rather than relying on the total score alone. Furthermore, with the replay tool, the users can perform more thorough diagnostics of the agent's behaviour, or better understand its benchmark scoring by visually inspecting the exact grid state and action taken at each timestep of a run. This is particularly valuable for detecting behaviour that is not directly captured by the metrics themselves, such as the reward hacking behaviour discussed in Chapter \ref{ch:experiments}.


\subsection{Configurable Components}
Both the scoring weights and the scenarios are configurable components of the benchmark. They determine the objective and quality of the benchmark. This flexibility allows \benchmark to be tailored to different research questions and evaluation goals, ranging from safety-oriented to cost-oriented evaluations, but it also shifts the responsibility onto the user to configure these components in a meaningful way. The following two subsections discuss the weights and the scenario selection individually, highlighting how each shapes the outcome of a benchmark trial.

\subsubsection{Weights}
In the lower layer of the scoring system, the metrics and categories have configurable weights. The user sets these to emphasise or omit specific parts of the evaluation. This makes it possible for researchers to create benchmarks that are tailored to their objectives. They can create a safety-oriented, cost-oriented, or computational-efficiency-oriented benchmark or a mix of them, which creates different ranking possibilities. 

Another use-case is shown in Chapter \ref{ch:experiments}, with two benchmarks focussing on grid operation and computational-efficiency. In the computational performance category, the do-nothing agent outperforms other agents, as it is used as the ideal reference point for minimal calculation complexity and time. It would not make sense to compare this agent directly with the greedy or \ac{ppo} agent. Thus, the configurable weights are used to create multiple benchmarking trials, which each analyse different aspects in power grid operation and evaluate different agent types.

\subsubsection{Scenario Selection} 
Another configurable part of a benchmark is its evaluation trials. In \benchmark, they are called scenarios, each covering one power grid, profile, grid modifications and specific indices in the profile. Together with the weights, they form a specific benchmark. The weights and scenarios define the objective of the benchmark and what it should evaluate. However, the scenarios also set the complexity and quality of the evaluations. Depending on the chosen scenarios, the coverage and difficulty can differ. The objective of the user might be diverse scenario or stress condition testing, or the users want to research a specific type of power grid and choose only that type as scenario. 

The importance in the design of \benchmark is that multiple scenarios are run in one benchmark trial. This ensures that the agents are tested on a high coverage dataset, which reveals whether the agent has strengths or weaknesses for specific scenarios. With only a single scenario, it is not possible to conduct a thorough analysis of an agent's capabilities in power grids.

However, this configurability also makes the resulting benchmark quality dependent on the user. If the user chooses poor configurations, then the benchmark might produce misleading conclusions about the agent's performances. Therefore, the selection of scenarios needs to be implemented with an objective in mind, and this objective should be pursued through multiple scenarios that have a high coverage of the possible cases.





\section{Evaluation of Benchmark Principles}\label{sec:evaluation-benchmark-principles}
Having established the principles of a well-designed benchmark in Section \ref{ssec:principles-benchmark}, this section evaluates the benchmark presented in this work against these criteria. Each principle is discussed individually, considering both the aspects that were successfully addressed and those that were only partially fulfilled.

\paragraph{Definition of Objective.}
As for clear objectives, \benchmark offers the user many choices, which define the objectives of the evaluation. For example, the user must decide which metrics and categories are important for evaluation and which scenarios should be covered to benchmark the agents. The benchmark then focusses on these cases and evaluates only these settings. The metrics and categories chosen for the trials cover important parts of optimal grid operation, considering grid stability and cost factors. These aspects follow the broader research objective of improving agents for power grid operation. However, this flexibility also means that \benchmark does not enforce or suggest a specific objective by default. Users unfamiliar with power grid operation may struggle to select a meaningful combination of metrics, categories and scenarios, which could lead to an evaluation that does not reflect a clearly defined or well-motivated objective.

\paragraph{Reproducibility.}
For reproducibility, the benchmark does not use any seeds for evaluation because the internal benchmark processes always run deterministically. The benchmark always performs the same tasks for the same input scenarios. After the agent runs through the different trials and the actions are saved internally, the evaluation process is always identical when the same run values are loaded, and identical scores are calculated. However, this only guarantees the reproducibility of the internal benchmark processes, not the reproducibility of the agent's behaviour itself. If an agent uses a stochastic policy or otherwise non-deterministic decision-making process, its actions, and therefore its resulting scores, may differ between runs of the same scenario, even though the benchmark process itself is fully deterministic. 
Additionally, stochasticity can be introduced for learned agents during training. For example, this can be achieved through random weight initialisation or non-deterministic exploration. This means that even a deterministic evaluation policy may produce different results for instances of the same agent that have been trained in different ways.
This limitation is discussed further in Section \ref{sec:discussion-limitations}.

\paragraph{Comparability.}
As seen in the experiments in Chapter \ref{ch:experiments}, \benchmark offers a structured comparability between multiple agents. It not only shows the best overall agent, but also the strengths and weaknesses among other agents and between different trials. Thus, it does not simply answer the question of which agent is the best. Instead, it provides an analysis of different cases and metrics inside power grids and a ranking over these features. However, in case of the \ac{ppo} agent's reward hacking, shown in the experiments, this was not captured by the metric system. The replay function made it possible to see more information on each agent, but it is not something measurable or comparable with the other agents. It merely shows whether the agent chooses questionable actions that are not optimal in power grid operation. 

\paragraph{Transparency.}
Through this work, the documentation\footnote{\documentation} and the evaluation results being stored in a human-readable JSON format (as presented in Section \ref{sec:int-processes}), the results are interpretable for the user. With the further visualisation tools presented in Section \ref{sec:result-format}, the user gains even more insight into the agent's performance during the benchmark experiments.

\paragraph{Relevance.}
As explained in Section \ref{sec:implementation-scenarios}, the power grid data for grid topology and profiles is provided by IEEE test cases and SimBench data. Both sources offer realistic power grid data for benchmarking power grid systems and cover different types of power grids. However, as discussed in Section \ref{sec:exp-analysis-interpretation}, some of these test cases contain physically implausible parameter values, limiting how closely the benchmark reflects real-world grid behaviour in these cases.

\paragraph{Representativeness.}
The quality of the benchmarking solely depends on the settings and the diversity of scenarios chosen by the user. This determines the representativeness, depending on the coverage of evaluated edge cases. The more scenarios are selected, especially with different grid modifications, the better the coverage is. However, when working in teams or in a competition, a set of diverse scenarios with high coverage and specified evaluation criteria can be predefined, allowing multiple agents to be compared and enabling a diverse benchmark set for research on agent improvement.





\section{Existing Benchmark Comparison}\label{sec:existing-benchmark-comparison}
Earlier in this work, several existing benchmarking tools for \ac{rl} agents in power grid operation were introduced, namely PowerGym, \ac{l2rpn} and RL2Grid, and summarised in Table \ref{tab:power-grid-benchmarks}. This section revisits these benchmarks and compares them directly against \benchmark, highlighting differences in scoring design, scenario diversity, safety handling and result presentation.

\paragraph{PowerGym.}
As discussed in Section \ref{ssec:power-benchmark}, PowerGym offers a benchmark on Volt-Var control systems in distribution grids. It uses the IEEE test cases for each evaluation and makes these cases modifiable like in \benchmark. As the objective of this benchmark is narrower, it makes the final scoring easier to interpret. The flexibility of \benchmark enables a broader evaluation, but at the cost of requiring additional configuration choices that directly influence the final ranking. However, the benchmarking feedback is measured as the cumulative reward, while the reward is composed of multiple grid metrics. This is a structural difference from \benchmark, with the feedback of PowerGym being single-dimensional. The experiments in Chapter \ref{ch:experiments} showed how crucial it is to analyse agents at multiple levels of depth to pinpoint unexpected problems, even when the overall score appears strong. Evaluating additional independent metrics or diagnostic tools can expose such behaviour.

\paragraph{L2RPN.}
The \ac{l2rpn} challenge provided realistic power-grid scenarios, with hidden evaluation scenarios, which are executed on a central server. For each submitted agent, it computed a standardised competition score, which was used for a leaderboard. With \benchmark, there is no evaluation server or leaderboard, and the flexibility in configuration makes it possible for two research groups to select completely different evaluation bases. The scores created by both research groups would not be comparable between the groups, due to differences in scenarios or score composition. Furthermore, due to the number of participants in \ac{l2rpn}, its design received substantially more external validation compared to the newly presented benchmark of this work. Additionally, the scoring also has a structural difference. In \ac{l2rpn}, it is relative to the do-nothing agent as a baseline, while in this work the do-nothing agent is used more as a selection model for choosing the scenarios. It is only used in the computational performance as a scoring reference. The advantage of the \ac{l2rpn} scoring system is, that each scenario is scored similarly, no matter how complex and difficult it is. If in a scenario it is harder to solve grid stability, then the do-nothing agent would also perform worse and the evaluated agents require similar effort to achieve a score above 0, as with easier scenarios. In \benchmark, the scoring also depends on the scenario difficulty, as the experiments in Chapter \ref{ch:experiments} showed. When a scenario was impossible to operate, as the first iteration already reached non-convergence, then all the agents received the lowest score. However, the scoring in \ac{l2rpn} was not fully transparent. While \benchmark uses a linear scoring system with worst and best values, in \ac{l2rpn} it remains unclear exactly how values between the fixed score points of -100, 0, 80, and 100 are calculated. Furthermore, this score is again a cumulative score, focused on competitive ranking for the leaderboard. This takes information away from single runs and makes diagnostics impossible. Moreover, the evaluation in \ac{l2rpn} is conducted only on one core grid environment, instead of offering a variety of different grids to test. Even though the experiments shown in Chapter \ref{ch:experiments} only covered two grids, \benchmark offers other grids from the IEEE test cases and SimBench data.

\paragraph{RL2Grid.}
RL2Grid builds on \ac{l2rpn} and comes closest to the benchmark of this work. It improves \ac{l2rpn} by adding more grid environments using Grid2Op systems, synthetic ChroniX2Grid profiles, expert-informed heuristics that emulate human operator behaviour and \ac{cmdp}-based safety constraints. The safety constraints offer the strongest advantage over \benchmark, as it improves safety by avoiding load shedding and islanding and thermal-limit violations. While load shedding and thermal limits are covered in the metrics of this work, load islanding is not part of the evaluation. The case study in Section \ref{ssec:case-studies} demonstrates that safety-critical behaviour may not always be adequately represented by metrics, thus leading the \ac{ppo} agent to choose actions that caused load shedding and islanding. This is avoidable using a constraint system, just like in RL2Grid. Additionally, RL2Grid includes an expert-informed heuristic, which suppresses unnecessary actions when the grid is already stable enough, so that the heuristic can incrementally restore the original topology. This covers a stronger integration of \ac{tso} operational knowledge into the benchmark process, which \benchmark does not possess. This is also an advantage for the agent's training and testing, as it reduces the task to challenging parts, where the agent needs to stabilise the grid. In terms of grid diversity, both tools offer a solid amount of usable scenario data, with the Grid2Op and ChroniX2Grid data in RL2Grid and the IEEE test cases and SimBench in \benchmark. Additionally, some of RL2Grid's larger grids include battery storage units, a feature not currently modelled in \benchmark's grid data.
As for the evaluation metrics, \benchmark offers more dimensions compared to the survival metrics and the three operational metrics (line margin, overloads and topology changes) in RL2Grid. This makes comparison over different agent strategies more limited for RL2Grid. Furthermore, the category of computational performance is not included in RL2Grid. This category is crucial in real-time control, as a policy that cannot meet operational timing constraints may be practically unusable. Another advantage of the benchmark in this work is the different formats in which results can be displayed. \benchmark offers a result-exploration layer, which enables in-depth analysis of the results and makes it easier to pinpoint outliers. The replay functionality offers fault detection for the agent's scoring by displaying its actions. Another disadvantage of RL2Grid is its use of random starting points or random line disconnections, reducing exact reproducibility, while \benchmark uses fixed configurations for the whole trial.





\section{Learnings from the Experiments}\label{sec:experiment-learnings}
The experiments conducted in Chapter \ref{ch:experiments} did not only serve to demonstrate the capabilities of \benchmark, but also provided broader insights into how \ac{rl} agents for power grid operation should be benchmarked in general. This section summarises four such learnings, regarding agent comparison, score aggregation, the role of baseline agents, and the influence of scenario configuration on the resulting conclusions.

\paragraph{Universally Strongest Agent.}
When comparing the three agents in the experiments, the results showed, that there is no single agent, that was universally strongest. Some agents excel in some scenarios, but fail in others. This was visible with the \ac{ppo} and do-nothing agents, that often competed for first place in the scenario scores. In some runs one agent was stronger, in others the other was stronger. It was never just one agent beating all of them. Agent quality is scenario-dependent and a benchmark has to illustrate this, to show strengths and weaknesses.

\paragraph{Score Aggregation Conceals Information.}
When looking at the \ac{ppo} agent, it overall achieved the best scores over the whole experiment. However, when taking a closer look, the agent seemed to have found a way to achieve high rewards, but still operate in unsuitable ways for power grid operation. This was discussed in Chapter \ref{ch:experiments} as a case of reward hacking, which made the \ac{ppo} agent achieve high scores but exhibit low usability in real-world operations. In summary, a high benchmark performance does not automatically imply operationally desirable behaviour. The actions of each agent should be analysed in depth, before deciding how well an agent performs.

\paragraph{Baselines Provide Context.}
Another finding from the experiments was the usefulness of the baseline agents. Even though the two baseline agents used simpler strategies, compared to the trained \ac{ppo} agent, they were still useful for the experiment interpretations. For example, in some scenarios the do-nothing agent survived surprisingly long, as shown in Section \ref{ssec:case-studies}. In contrast, the greedy agent acts as the opposite of the do-nothing agent, by intervening frequently without any long-term robustness in many scenarios. Both baselines provide information in each scenario, not only for ranking against other agents, but also for interpreting the scenario's difficulty. The worse the baselines perform, the harder it is for other, more complex agents to maintain grid stability.

\paragraph{Configured Scenarios Affect Conclusions.}
As discussed in Section \ref{sec:discussion-benchmark-design}, the scenarios chosen by the user strongly affect the conclusions of the benchmarking. Each of the grid modification experiments in Chapter \ref{ch:experiments} demonstrates that the outcomes differ. When using only a part of these experiments, the conclusions are more limited and result in a different scoring and thus ranking of each agent. Through the experiments, it is clear that one scenario or just one run is not sufficient for a deep understanding of each agent's performance. A diverse set of scenarios is essential for a stronger benchmark quality and a better analysis of each agent. 




\section{Limitations}\label{sec:discussion-limitations}
While \benchmark addresses the research question posed in this thesis, the experiments conducted with it also revealed several limitations that should be considered when interpreting the presented results. This section discusses five such limitations, concerning the selection of agents and grids used in the experiments, the quality of the underlying grid data, the handling of stochastic agent behaviour, and the computational evaluation environment.

\paragraph{Limited Agent Selection.}
During the experiments, only three agents were evaluated. Additionally, only one of these agents has a learned \ac{rl} architecture. For this reason, this work did not fully validate \benchmark across different \ac{rl} architectures, such as those discussed in Section \ref{ssec:example-agents}. A broader comparison across value-based, policy-based and hybrid architectures, as discussed in Section \ref{ssec:future-agents}, would be necessary to assess whether the observed differences in benchmark scores generalise across agent types.

\paragraph{Limited Grid Selection.}
Another limitation in the experiments is the lack of variety in different grid topologies. Only two IEEE test cases were evaluated, but the benchmark supports more grids. The benchmark focussed more on grid modification and finding challenging scenarios, but it lacks testing on larger transmission grids or real-world grids. This limits the ability to draw conclusions about how \benchmark scales to larger, more interconnected grids, where topology actions may have different effects on grid stability.

\paragraph{Grid Quality.}
As discussed in Section \ref{sec:exp-analysis-interpretation}, the 14-bus grid from the experiments has some physically implausible values for the grid edges, such as the \textit{vk\_percent} values in the transformers or the \textit{max\_i\_ka} values in the transmission lines. For this reason, some observed instability may originate partly from the grid model rather than the agent performance. This limits the interpretation of some experimental results. This expands the lesson from Sections \ref{sec:discussion-benchmark-design}, \ref{sec:evaluation-benchmark-principles} and \ref{sec:experiment-learnings} about the benchmark quality. The benchmark quality is bounded by the selection of the scenarios \textbf{and} the quality of the scenario model.

\paragraph{Stochastic Evaluation.}
The benchmark scenarios and scoring procedure are deterministic for a fixed configuration. However, individual agents may introduce stochastic action selection, depending on their implementation. For this reason, running each of the scenario runs only once would not be practical in case of stochastic decisions inside the evaluated agent. Instead, each run needs to be repeated multiple times to calculate an average score for each metric score of a run. Furthermore, the user should be able to set one or multiple seeds, to have a fixed testing environment for these types of agents. This seed needs to be used by the agent, so that the stochastic decision-making is reproducible. In addition to stochasticity at evaluation time, learned agents may also be affected by stochasticity introduced during training, such as random parameter initialisation or exploration noise. Consequently, two agents trained using the same algorithm and hyperparameters but different training seeds may produce different policies and thus achieve different benchmark scores. Therefore, a thorough evaluation should consider training multiple instances of the same agent with different training seeds to determine how much of the measured performance is attributable to the algorithm itself rather than a specific training run.

\paragraph{Computational Benchmark Environment.}
As the experiment testing computational efficiency was only run on one machine, the values are hardware-specific. To produce a stronger and more general evaluation of the two tested agents, it is necessary to test the computational efficiency on multiple machines. Additionally, computational performance measurements are inherently variable as they can be affected by external factors, such as background processes or system load, at the time of measurement. Therefore, a single run is not sufficient to obtain a reliable estimate, similar to the handling of stochastic agent behaviour discussed above. Multiple repetitions of each measurement are necessary to obtain stable and representative values. Furthermore, the current evaluation only measures an agent's overall inference time, without profiling or tracing its internal computations. Consequently, it is impossible to determine why one agent is faster or slower than another or which parts of its decision-making process are computationally expensive. Agents can be implemented with varying degrees of efficiency, for example through optimised code or better use of hardware acceleration, but this simple timing measurement alone cannot capture or distinguish such implementation-specific differences. Future work could address this by integrating dedicated profiling tools, allowing a more detailed breakdown of an agent's computation time across its decision-making pipeline.